% =============================================================
%  开局说明
% =============================================================
\Chapter{开局说明：这套东西怎么玩}

\begin{unitinfo}
先说清楚：\textbf{这不是练习册。}里面一道“口算天天练”都没有。\\
它是一份\textbf{陪小豆角闯关的说明书}——豆角爸爸照着念、照着做，
二十分钟一到，收工，谁也不许加练。
\end{unitinfo}

\subsection{出场人物}

\begin{center}
\begin{tabularx}{\linewidth}{@{}p{2.6cm}X@{}}
\toprule
\textbf{小豆角} & 二年级，本书唯一的服务对象。强项是记性好，弱项是缺乏逻辑思维且在数学上长期得不到正向反馈，导致
对数学约错越怕，越怕越错，对数学产生恐惧。 \\
\textbf{豆角爸爸} & 本书作者、编辑、校对，以及\textbf{头号翻车选手}。
心算经常错，量东西经常量错，剪窗花像剪怪兽。\\
\textbf{棒子弟} & 那包小棒。一百根，捆起来的时候会“变身”——
十根一捆，就变成了“一个十”。 \\
\textbf{啪嗒哥} & 五档计数器。往上推珠子会“啪嗒”一声，
\textbf{借位的时候它最忙}，连着啪嗒两下。 \\
\textbf{剪刀姐} & 专门负责对称。她对“折痕有没有对齐”这件事非常挑剔。 \\
\bottomrule
\end{tabularx}
\end{center}

\subsection{闯关地图}

一共八关，加一个隐藏关，最后是终局。\textbf{每一关都对着课本的一个单元}，
所以学校讲到哪儿，豆角爸爸就跳到哪一关。

\begin{center}
\begin{tikzpicture}[font=\scriptsize,
  box/.style={draw=mainc,rounded corners=1.2mm,fill=mainc!6,inner sep=3pt,
              minimum width=3.45cm,minimum height=0.95cm,align=center,text width=3.4cm}]
  \node[box] (n1) at (0,0) {\textbf{1}\ 一串小棒走天下};
  \node[box] (n2) at (3.95,0) {\textbf{2}\ 尺子、绳子和一双大手};
  \node[box] (n3) at (7.9,0) {\textbf{3}\ 把加法压扁};
  \node[box] (n4) at (11.85,0) {\textbf{4}\ 编顺口溜的人};
  \node[box] (n5) at (11.85,-1.35) {\textbf{5}\ 分东西，不许打架};
  \node[box] (n6) at (7.9,-1.35) {\textbf{6}\ 剪一刀，动一下};
  \node[box] (n7) at (3.95,-1.35) {\textbf{7}\ 八九七十二的来历};
  \node[box] (n8) at (0,-1.35) {\textbf{8}\ 算账的本事};
  \node[box] (n9) at (0,-2.7) {\textbf{隐藏}\ 画地图，花真钱};
  \node[box] (n10) at (3.95,-2.7) {\textbf{终局}\ 考前三十天};
  \foreach \a/\b in {n1/n2,n2/n3,n3/n4,n5/n6,n6/n7,n7/n8}{
    \draw[-{Stealth},mainc,thick] (\a) -- (\b);}
  \draw[-{Stealth},mainc,thick] (n4) -- (n5);
  \draw[-{Stealth},mainc,thick] (n8) -- (n9);
  \draw[-{Stealth},mainc,thick] (n9) -- (n10);
\end{tikzpicture}
\end{center}

\noindent{\small\color{tipc}
对照：第 1 关＝课本第一单元（100 以内加与减）· 第 2 关＝测量 · 第 3 关＝数一数与乘法 ·
第 4 关＝乘法口诀（一）· 第 5 关＝分一分与除法 · 第 6 关＝图形的运动 ·
第 7 关＝乘法口诀（二）· 第 8 关＝乘除法的应用 · 隐藏关＝两次综合实践 · 终局＝总复习}

\subsection{每一关里有什么}

每关的结构都一样。重点是下面这三件“玩的东西”，
它们不是点缀，\textbf{是让小豆角愿意坐下来的唯一理由}。

\begin{blunder}
\textbf{豆角爸爸的翻车现场}。每关开头，爸爸先自曝一个真实糗事：
量错床单、剪出怪兽、算错找零。
\textbf{作用}：让小豆角知道“错”不是丢人的事，爸爸错得比小豆角多。
\end{blunder}

\begin{gamble}{课前协商：这一页最要紧}
每关开头都有一个\textbf{课前协商}框：\textbf{今天比什么}由这本书来定，
\textbf{赢了拿什么、输了认什么}留空——\textbf{上课前两分钟谈好，当场手写填上去}。\\[4pt]
根据豆角现阶段最想要的东西设计正向激励。\\[5pt]
\textbf{三条铁律}（很重要，别破坏）：
\begin{enumerate}[label=\arabic*、]
  \item \textbf{豆角爸爸一定要真的输一次。}全程赢的爸爸，小豆角不会想再跟爸爸玩了。
  \item \textbf{赌注要小、要当场兑现。}小到不用心疼，近到当天能兑现；
        别赌“考好了给小豆角买xx”。
  \item \textbf{先让小豆角开口提。}她自己提的赌注，她才认账；爸爸只负责把它压小一点。
\end{enumerate}
\end{gamble}

\begin{bonus}
每关还有一个\textbf{跟分数完全没关系}的野任务：电梯挑战、身高排行榜、
把口诀唱成 rap、藏宝图、当一天家庭采购员。\\
\textbf{作用}：把数学塞进小豆角本来就爱干的事里。做了不加分，不做也不骂。
\end{bonus}

\subsection{二十分钟怎么切}

说实话，真正“讲”的时间只有两分钟。剩下全是小豆角在动手。
这不是偷懒，这是小豆角唯一能吸收的方式。

\begin{center}
\begin{tabularx}{\linewidth}{@{}p{1.4cm}p{2.3cm}X@{}}
\toprule
\textbf{时间} & \textbf{干什么} & \textbf{爸爸干什么} \\
\midrule
3 分钟 & 抽卡口算 & 抽十张，小豆角报答案。爸爸只记对错，\textbf{当场不纠}。\\
2 分钟 & 小豆角自己读题 & 爸爸只说一句“把数字圈出来”。\textbf{爸爸不念题}，一念小豆角就不看题了。\\
8 分钟 & \textbf{上手摆} & 全场最重要的八分钟。爸爸发东西，小豆角摆，爸爸不可以说话。\\
5 分钟 & 落到纸上 & 小豆角写算式、竖式或口诀。\textbf{一个字爸爸都不代笔}。\\
2 分钟 & 收尾 & 三句：今天学了啥？哪步最容易翻车？明天还来吗？\\
\bottomrule
\end{tabularx}
\end{center}

\subsection{三种开法，挑一种}

\begin{center}
\begin{tabularx}{\linewidth}{@{}p{2.6cm}p{2.4cm}X@{}}
\toprule
\textbf{开法} & \textbf{节奏} & \textbf{适合什么情况} \\
\midrule
\textbf{跟着学校走} & 一周 2 次 & 周二、周四晚上各一次，
学校讲完哪个单元，爸爸就翻到那一关。一学期 41 次，刚好放到期末。\\
\textbf{假期先扫一遍} & 一周 5 次 & 寒暑假集中过，开学前把整册摸完。
开学后上课就等于复习，小豆角会很轻松。\\
\textbf{哪儿着火灭哪儿} & 一周 3 次 & 竖式老错、口诀背不下来，
就只打那一关。先做关末的抽查，找出病根，再往回补。\\
\bottomrule
\end{tabularx}
\end{center}

\noindent 不管选哪种，这三条别动：
\textbf{① 一次不超过 25 分钟；② 一次只推进一个点；③ 结束前让小豆角自己复述一遍。}

\subsection{爸爸的约定}

\begin{enumerate}
  \item \textbf{二十分钟到点就停。}没做完也停。宁可明天再来，别今天把小豆角磨烦。
  \item \textbf{不抢答。}小豆角停五秒爸爸才开口。五秒是小豆角把话组织出来的正常时间。
  \item \textbf{新的东西，第一次必须上手。}小棒也好纸片也好手指也好，
        先摆出来，摆完才许写算式。
  \item \textbf{小豆角做错了，爸爸先问“小豆角是怎么想的”。}不先说“错了”。
        让小豆角把错路走完，爸爸才知道断在哪。
  \item \textbf{一次只教一个点。}今天搞“进位”，就别顺嘴讲“退位”。
        讲了等于两个都没讲。
  \item \textbf{“这么简单都不会”这句话，一次都不说。}换成“这步确实容易忘，再看一遍”。
  \item \textbf{爸爸也不许掏手机算。}宁可一起重新摆一遍。
  \item \textbf{允许小豆角比爸爸慢。}小豆角摆十五根小棒花四十秒，正常。
  \item \textbf{每次留一个好印象。}收尾必须指着一件小豆角做对的事说清楚：
        “今天小豆角自己想到要捆起来，这个很关键。”
  \item \textbf{课前谈好的当场兑现。}输了就认，不许拖到明天。
\end{enumerate}

\noindent{\small\color{tipc}正文里出现\,\waterdrop\,的地方，是给爸爸的暗号：
“这里可以让一步，放点水。”小豆角不用懂。}

\subsection{爸爸自己踩过的六个坑}

\begin{pitbox}
\begin{enumerate}[label=\arabic*、]
  \item \textbf{把规则直接说出口。}“个位满十向十位进一”，爸爸一说，小豆角抄下来，
        看着像会了。换一组数又不会了。\\
        \textbf{改法}：让小豆角自己把十根捆成一捆，然后问小豆角“这一捆从哪来的”。
        让小豆角自己说出“十根换一捆”——这句话是小豆角说的，才管用。
  \item \textbf{拿口诀顶替理解。}小豆角能背“六七四十二”，但不知道怎么画出来。
        题目一改成“画图表示”，交白卷。\\
        \textbf{改法}：每句新口诀，先在点子图或百数表上圈出那块地方，圈完再背。
  \item \textbf{管小豆角掰手指。}爸爸一开始也管，说“多大了还掰手指”。
        结果小豆角改成在桌子底下偷偷数，还多了一层紧张。\\
        \textbf{改法}：见下一节，掰手指有明确的时间表，到点自然撤。
  \item \textbf{用题量顶思考。}一口气一百道口算，小豆角刷完还是不会。\\
        \textbf{改法}：一种类型做透五道，比五十道管用。
  \item \textbf{边做边纠错。}小豆角写一步爸爸说一次“不对”，最后小豆角干脆停手等爸爸报答案。\\
        \textbf{改法}：整道题做完，再一起回头看。
  \item \textbf{盯着分数。}二年级上册真正的目标是“能说清为什么”。
        说不清却做对，是运气；说得清但算错，是进步。
\end{enumerate}
\end{pitbox}

\subsection{掰手指这事儿，说清楚}

不管，但要有退出时间表。

\begin{center}
\begin{tabularx}{\linewidth}{@{}p{3.6cm}X@{}}
\toprule
\textbf{什么时候} & \textbf{允许小豆角用手吗} \\
\midrule
第 1 关 & 随便掰。\textbf{别催}。这阶段目标是“懂”，不是“快”。\\
第 2--4 关 & 20 以内还能掰；两位数的进位加、退位减
\textbf{必须用小棒或竖式}。手指不够用，正好自然过渡。\\
第 5--8 关 & 口诀必须脱手，直接报数；除法可以心里想口诀。\\
终局 & 100 以内加减，允许小豆角在心里想竖式，不再动手。\\
\bottomrule
\end{tabularx}
\end{center}

\noindent \textbf{怎么判断该不该管}，看两件事就够：

\begin{itemize}
  \item \textbf{又慢又准} $\to$ 正常，别管，小豆角在认真算。
  \item \textbf{又快又错} $\to$ 小豆角在猜，必须退回“摆”。
  \item \textbf{又慢又错} $\to$ 某一步根本没懂，退回上一关重来。
\end{itemize}

\subsection{错题本：一次只写一句话}

不要抄整道题。抄题小豆角会烦，抄完也不会看。

\begin{itemize}
  \item 这样写：\ “48+25，个位 8+5=13，爸爸把进的 1 忘了加到十位。”
  \item 不要这样写：\ “粗心”“算错了”“不认真”——一条都别写，写了等于没写。
\end{itemize}

\noindent 每周日花五分钟，让小豆角读一遍自己写的错因，然后口头重做那三道题。做对就划掉。

\textbf{同一个错因出现三次，停课，回去补那一关。}
这是整本书最重要的一条判断依据，比任何分数都准。

\subsection{战利品：星星，和赢爸爸一次}

二年级小豆角对分数没感觉，对贴纸和“赢爸爸一次”很有感觉。

\begin{itemize}
  \item 一次课顺顺当当上完（不哭、不跑、没超时）$\to$ \starx\ 一颗。
  \item 一次抽查全对 $\to$ \starx\ 一颗。
  \item \textbf{小豆角自己发现并改了一个错} $\to$ \starx\starx\ 两颗。
        \textbf{这条最重要}，爸爸在刻意鼓励小豆角自查，这个比做对更有价值。
  \item 攒够 20 颗，换一次小豆角自己点单的东西。
\end{itemize}

\noindent \textbf{底线}：不用“考一百分”、“考第一”当条件。

\subsection{每次开打前，三分钟准备}

\begin{enumerate}
  \item \cb 课本翻到这一关标的页码了吗？
  \item \cb “先备货”框里的东西都拿到了吗？缺一样就先去找，别用“凑合”代替。
  \item \cb 桌子清空了吗？只留课本、笔、材料。玩具和平板收走。
  \item \cb 计时器设 20 分钟，放在小豆角看得见的地方。
  \item \cb 爸爸自己把“照这样问小豆角”看过一遍了吗？\textbf{提前看，别现场编。}
\end{enumerate}

\noindent 五条都打钩再开始。\textbf{准备不足就改天，别硬上}——
硬上浪费时间。
